% BEGIN GENERATED ARTIFACT: def:propositional-equivalence-law
\begin{definitionbox}{Definition (Propositional Equivalence Law)}
\begin{definition}[Propositional Equivalence Law]
\label{def:propositional-equivalence-law}
A propositional equivalence law is a valid schema of the form
\[
\varphi\equiv\psi,
\]
with \(\varphi,\psi\in\WFF\). Each instance of the schema says that the two formulas have the same truth value under every truth assignment.
\end{definition}
\end{definitionbox}

\begin{remark*}[Standard quantified statement]
\[
\operatorname{PropositionalEquivalenceLaw}(\varphi,\psi)
\Longleftrightarrow
\varphi,\psi\in\WFF
\land
\varphi\equiv\psi.
\]
\end{remark*}

\begin{remark*}[Definition predicate reading]
\[
\operatorname{PropositionalEquivalenceLaw}(\varphi,\psi)
\]
means that \(\varphi\equiv\psi\) is a valid logical-equivalence schema.
\end{remark*}

\begin{remark*}[Interpretation]
Equivalence laws are the algebraic rewrite rules of propositional logic. They allow formulas to be transformed without changing truth conditions.
\end{remark*}

\begin{dependencies}
\begin{itemize}
  \item \hyperref[def:logical-equivalence-propositional-logic]{Logical Equivalence}
  \item \hyperref[def:well-formed-formula]{Well-Formed Formula}
\end{itemize}
\end{dependencies}
% END GENERATED ARTIFACT: def:propositional-equivalence-law
% BEGIN GENERATED ARTIFACT: thm:logical-equivalence-equivalence-relation
\begin{theorembox}{Theorem (Logical Equivalence is an Equivalence Relation)}
\begin{theorem}[Logical Equivalence is an Equivalence Relation]
\label{thm:logical-equivalence-equivalence-relation}
Logical equivalence is reflexive, symmetric, and transitive on \(\WFF\):
\[
\varphi\equiv\varphi,
\qquad
\varphi\equiv\psi\Longrightarrow\psi\equiv\varphi,
\]
and
\[
(\varphi\equiv\psi\land\psi\equiv\chi)
\Longrightarrow
\varphi\equiv\chi.
\]

\noindent\hyperref[prf:logical-equivalence-equivalence-relation]{Go to proof.}
\end{theorem}
\end{theorembox}

\begin{remark*}[Standard quantified statement]
\[
\begin{aligned}
&(\forall\varphi\in\WFF)(\varphi\equiv\varphi),\\
&(\forall\varphi,\psi\in\WFF)(\varphi\equiv\psi\to\psi\equiv\varphi),\\
&(\forall\varphi,\psi,\chi\in\WFF)
((\varphi\equiv\psi\land\psi\equiv\chi)\to\varphi\equiv\chi).
\end{aligned}
\]
\end{remark*}

\begin{remark*}[Interpretation]
Logical equivalence behaves like equality at the semantic level.
\end{remark*}

\begin{dependencies}
\begin{itemize}
  \item \hyperref[def:logical-equivalence-propositional-logic]{Logical Equivalence}
\end{itemize}
\end{dependencies}
% END GENERATED ARTIFACT: thm:logical-equivalence-equivalence-relation
% BEGIN GENERATED ARTIFACT: thm:basic-boolean-equivalence-laws-propositional-logic
\begin{theorembox}{Theorem (Basic Boolean Equivalence Laws)}
\begin{theorem}[Basic Boolean Equivalence Laws]
\label{thm:basic-boolean-equivalence-laws-propositional-logic}
For all formulas \(\varphi,\psi,\chi\in\WFF\), the following logical equivalences hold.

Commutativity:
\[
\varphi\land\psi\equiv\psi\land\varphi,
\qquad
\varphi\lor\psi\equiv\psi\lor\varphi.
\]

Associativity:
\[
(\varphi\land\psi)\land\chi
\equiv
\varphi\land(\psi\land\chi),
\qquad
(\varphi\lor\psi)\lor\chi
\equiv
\varphi\lor(\psi\lor\chi).
\]

Idempotence:
\[
\varphi\land\varphi\equiv\varphi,
\qquad
\varphi\lor\varphi\equiv\varphi.
\]

Double negation:
\[
\neg\neg\varphi\equiv\varphi.
\]

\noindent\hyperref[prf:basic-boolean-equivalence-laws-propositional-logic]{Go to proof.}
\end{theorem}
\end{theorembox}

\begin{remark*}[Standard quantified statement]
\[
(\forall\varphi,\psi,\chi\in\WFF)
\left[
\begin{gathered}
\varphi\land\psi\equiv\psi\land\varphi,\quad
\varphi\lor\psi\equiv\psi\lor\varphi,\\
(\varphi\land\psi)\land\chi\equiv\varphi\land(\psi\land\chi),\quad
(\varphi\lor\psi)\lor\chi\equiv\varphi\lor(\psi\lor\chi),\\
\varphi\land\varphi\equiv\varphi,\quad
\varphi\lor\varphi\equiv\varphi,\quad
\neg\neg\varphi\equiv\varphi
\end{gathered}
\right].
\]
\end{remark*}

\begin{remark*}[Interpretation]
These are the basic algebraic laws for conjunction, disjunction, and negation.
\end{remark*}

\begin{dependencies}
\begin{itemize}
  \item \hyperref[def:logical-equivalence-propositional-logic]{Logical Equivalence}
  \item \hyperref[thm:unique-extension-truth-assignment-propositional-logic]{Unique Extension of a Truth Assignment}
\end{itemize}
\end{dependencies}
% END GENERATED ARTIFACT: thm:basic-boolean-equivalence-laws-propositional-logic
% BEGIN GENERATED ARTIFACT: thm:identity-domination-laws-propositional-logic
\begin{theorembox}{Theorem (Identity and Domination Laws)}
\begin{theorem}[Identity and Domination Laws]
\label{thm:identity-domination-laws-propositional-logic}
Let \(\top\) be any tautological formula and let \(\bot\) be any contradictory formula. Then for every \(\varphi\in\WFF\):
\[
\varphi\land\top\equiv\varphi,
\qquad
\varphi\lor\bot\equiv\varphi,
\]
and
\[
\varphi\land\bot\equiv\bot,
\qquad
\varphi\lor\top\equiv\top.
\]

\noindent\hyperref[prf:identity-domination-laws-propositional-logic]{Go to proof.}
\end{theorem}
\end{theorembox}

\begin{remark*}[Standard quantified statement]
\[
\begin{aligned}
&(\operatorname{Tautology}(\top)\land \operatorname{Contradiction}(\bot))
\to
(\forall\varphi\in\WFF)\\
&\quad
\left[
(\varphi\land\top\equiv\varphi)
\land
(\varphi\lor\bot\equiv\varphi)
\land
(\varphi\land\bot\equiv\bot)
\land
(\varphi\lor\top\equiv\top)
\right].
\end{aligned}
\]
\end{remark*}

\begin{remark*}[Interpretation]
Tautologies act like true constants; contradictions act like false constants.
\end{remark*}

\begin{dependencies}
\begin{itemize}
  \item \hyperref[def:tautology-propositional-logic]{Tautology}
  \item \hyperref[def:contradiction-propositional-logic]{Contradiction}
  \item \hyperref[def:logical-equivalence-propositional-logic]{Logical Equivalence}
\end{itemize}
\end{dependencies}
% END GENERATED ARTIFACT: thm:identity-domination-laws-propositional-logic
% BEGIN GENERATED ARTIFACT: thm:de-morgan-laws-propositional-logic
\begin{theorembox}{Theorem (De Morgan Laws for Propositional Logic)}
\begin{theorem}[De Morgan Laws for Propositional Logic]
\label{thm:de-morgan-laws-propositional-logic}
For all \(\varphi,\psi\in\WFF\),
\[
\neg(\varphi\land\psi)
\equiv
\neg\varphi\lor\neg\psi,
\]
and
\[
\neg(\varphi\lor\psi)
\equiv
\neg\varphi\land\neg\psi.
\]

\noindent\hyperref[prf:de-morgan-laws-propositional-logic]{Go to proof.}
\end{theorem}
\end{theorembox}

\begin{remark*}[Standard quantified statement]
\[
\begin{gathered}
(\forall\varphi,\psi\in\WFF)
\left[
\neg(\varphi\land\psi)
\equiv
\neg\varphi\lor\neg\psi
\right],\\
(\forall\varphi,\psi\in\WFF)
\left[
\neg(\varphi\lor\psi)
\equiv
\neg\varphi\land\neg\psi
\right].
\end{gathered}
\]
\end{remark*}

\begin{remark*}[Predicate reading]
\[
\underbrace{\neg(\varphi\land\psi)}_{\text{not both}}
\equiv
\underbrace{\neg\varphi\lor\neg\psi}_{\text{at least one fails}},
\qquad
\underbrace{\neg(\varphi\lor\psi)}_{\text{neither}}
\equiv
\underbrace{\neg\varphi\land\neg\psi}_{\text{both fail}}.
\]
\end{remark*}

\begin{remark*}[Interpretation]
Negation swaps conjunction with disjunction when pushed inward.
\end{remark*}

\begin{dependencies}
\begin{itemize}
  \item \hyperref[def:logical-equivalence-propositional-logic]{Logical Equivalence}
  \item \hyperref[thm:unique-extension-truth-assignment-propositional-logic]{Unique Extension of a Truth Assignment}
\end{itemize}
\end{dependencies}
% END GENERATED ARTIFACT: thm:de-morgan-laws-propositional-logic
% BEGIN GENERATED ARTIFACT: thm:distributive-laws-propositional-logic
\begin{theorembox}{Theorem (Distributive Laws for Propositional Logic)}
\begin{theorem}[Distributive Laws for Propositional Logic]
\label{thm:distributive-laws-propositional-logic}
For all \(\varphi,\psi,\chi\in\WFF\),
\[
\varphi\land(\psi\lor\chi)
\equiv
(\varphi\land\psi)\lor(\varphi\land\chi),
\]
and
\[
\varphi\lor(\psi\land\chi)
\equiv
(\varphi\lor\psi)\land(\varphi\lor\chi).
\]

\noindent\hyperref[prf:distributive-laws-propositional-logic]{Go to proof.}
\end{theorem}
\end{theorembox}

\begin{remark*}[Standard quantified statement]
\[
\begin{gathered}
(\forall\varphi,\psi,\chi\in\WFF)
\left[
\varphi\land(\psi\lor\chi)
\equiv
(\varphi\land\psi)\lor(\varphi\land\chi)
\right],\\
(\forall\varphi,\psi,\chi\in\WFF)
\left[
\varphi\lor(\psi\land\chi)
\equiv
(\varphi\lor\psi)\land(\varphi\lor\chi)
\right].
\end{gathered}
\]
\end{remark*}

\begin{remark*}[Interpretation]
Conjunction distributes over disjunction, and disjunction distributes over conjunction. These laws are the main algebraic engine behind CNF and DNF conversion.
\end{remark*}

\begin{dependencies}
\begin{itemize}
  \item \hyperref[def:logical-equivalence-propositional-logic]{Logical Equivalence}
  \item \hyperref[thm:unique-extension-truth-assignment-propositional-logic]{Unique Extension of a Truth Assignment}
\end{itemize}
\end{dependencies}
% END GENERATED ARTIFACT: thm:distributive-laws-propositional-logic
% BEGIN GENERATED ARTIFACT: thm:absorption-laws-propositional-logic
\begin{theorembox}{Theorem (Absorption Laws)}
\begin{theorem}[Absorption Laws]
\label{thm:absorption-laws-propositional-logic}
For all \(\varphi,\psi\in\WFF\),
\[
\varphi\land(\varphi\lor\psi)\equiv\varphi,
\]
and
\[
\varphi\lor(\varphi\land\psi)\equiv\varphi.
\]

\noindent\hyperref[prf:absorption-laws-propositional-logic]{Go to proof.}
\end{theorem}
\end{theorembox}

\begin{remark*}[Standard quantified statement]
\[
\begin{gathered}
(\forall\varphi,\psi\in\WFF)
\left[
\varphi\land(\varphi\lor\psi)\equiv\varphi
\right],\\
(\forall\varphi,\psi\in\WFF)
\left[
\varphi\lor(\varphi\land\psi)\equiv\varphi
\right].
\end{gathered}
\]
\end{remark*}

\begin{remark*}[Interpretation]
Once \(\varphi\) is already required, adding \(\varphi\lor\psi\) contributes no additional constraint. Dually, once \(\varphi\) is already allowed, adding \(\varphi\land\psi\) contributes no additional possibility.
\end{remark*}

\begin{dependencies}
\begin{itemize}
  \item \hyperref[def:logical-equivalence-propositional-logic]{Logical Equivalence}
  \item \hyperref[thm:basic-boolean-equivalence-laws-propositional-logic]{Basic Boolean Equivalence Laws}
  \item \hyperref[thm:distributive-laws-propositional-logic]{Distributive Laws for Propositional Logic}
\end{itemize}
\end{dependencies}
% END GENERATED ARTIFACT: thm:absorption-laws-propositional-logic
% BEGIN GENERATED ARTIFACT: thm:conditional-equivalence-laws-propositional-logic
\begin{theorembox}{Theorem (Conditional Equivalence Laws)}
\begin{theorem}[Conditional Equivalence Laws]
\label{thm:conditional-equivalence-laws-propositional-logic}
For all \(\varphi,\psi\in\WFF\),
\[
\varphi\to\psi
\equiv
\neg\varphi\lor\psi,
\]
and
\[
\neg(\varphi\to\psi)
\equiv
\varphi\land\neg\psi.
\]

\noindent\hyperref[prf:conditional-equivalence-laws-propositional-logic]{Go to proof.}
\end{theorem}
\end{theorembox}

\begin{remark*}[Standard quantified statement]
\[
\begin{gathered}
(\forall\varphi,\psi\in\WFF)
\left[
\varphi\to\psi
\equiv
\neg\varphi\lor\psi
\right],\\
(\forall\varphi,\psi\in\WFF)
\left[
\neg(\varphi\to\psi)
\equiv
\varphi\land\neg\psi
\right].
\end{gathered}
\]
\end{remark*}

\begin{remark*}[Predicate reading]
\[
\underbrace{\varphi\to\psi}_{\text{no counterexample}}
\equiv
\underbrace{\neg\varphi\lor\psi}_{\text{premise false or conclusion true}},
\qquad
\underbrace{\neg(\varphi\to\psi)}_{\text{counterexample}}
\equiv
\underbrace{\varphi\land\neg\psi}_{\text{premise true and conclusion false}}.
\]
\end{remark*}

\begin{remark*}[Interpretation]
A conditional fails exactly when the antecedent is true and the consequent is false.
\end{remark*}

\begin{dependencies}
\begin{itemize}
  \item \hyperref[def:logical-equivalence-propositional-logic]{Logical Equivalence}
  \item \hyperref[thm:unique-extension-truth-assignment-propositional-logic]{Unique Extension of a Truth Assignment}
\end{itemize}
\end{dependencies}
% END GENERATED ARTIFACT: thm:conditional-equivalence-laws-propositional-logic
% BEGIN GENERATED ARTIFACT: thm:biconditional-equivalence-laws-propositional-logic
\begin{theorembox}{Theorem (Biconditional Equivalence Laws)}
\begin{theorem}[Biconditional Equivalence Laws]
\label{thm:biconditional-equivalence-laws-propositional-logic}
For all \(\varphi,\psi\in\WFF\),
\[
\varphi\leftrightarrow\psi
\equiv
(\varphi\to\psi)\land(\psi\to\varphi),
\]
and
\[
\varphi\leftrightarrow\psi
\equiv
(\varphi\land\psi)\lor(\neg\varphi\land\neg\psi).
\]

\noindent\hyperref[prf:biconditional-equivalence-laws-propositional-logic]{Go to proof.}
\end{theorem}
\end{theorembox}

\begin{remark*}[Standard quantified statement]
\[
\begin{gathered}
(\forall\varphi,\psi\in\WFF)
\left[
\varphi\leftrightarrow\psi
\equiv
(\varphi\to\psi)\land(\psi\to\varphi)
\right],\\
(\forall\varphi,\psi\in\WFF)
\left[
\varphi\leftrightarrow\psi
\equiv
(\varphi\land\psi)\lor(\neg\varphi\land\neg\psi)
\right].
\end{gathered}
\]
\end{remark*}

\begin{remark*}[Predicate reading]
\[
\varphi\leftrightarrow\psi
\quad\Longleftrightarrow\quad
(\varphi\to\psi)\land(\psi\to\varphi)
\]
and also
\[
\varphi\leftrightarrow\psi
\quad\Longleftrightarrow\quad
(\varphi\land\psi)\lor(\neg\varphi\land\neg\psi).
\]
The biconditional says that the two formulas have the same truth value.
\end{remark*}

\begin{remark*}[Interpretation]
A biconditional says that two formulas have the same truth value: both true or both false.
\end{remark*}

\begin{dependencies}
\begin{itemize}
  \item \hyperref[def:logical-equivalence-propositional-logic]{Logical Equivalence}
  \item \hyperref[thm:conditional-equivalence-laws-propositional-logic]{Conditional Equivalence Laws}
\end{itemize}
\end{dependencies}
% END GENERATED ARTIFACT: thm:biconditional-equivalence-laws-propositional-logic
% BEGIN GENERATED ARTIFACT: def:formula-context-propositional-logic
\begin{definitionbox}{Definition (Formula Context)}
\begin{definition}[Formula Context]
\label{def:formula-context-propositional-logic}
A formula context \(C[-]\) is a well-formed formula with one distinguished placeholder occurrence \([-]\). For each well-formed formula \(\varphi\), the expression \(C[\varphi]\) denotes the formula obtained by replacing the placeholder with \(\varphi\).
\end{definition}
\end{definitionbox}

\begin{remark*}[Standard quantified statement]
\[
\operatorname{FormulaContext}(C[-])
\Longrightarrow
(\forall\varphi\in\WFF)(C[\varphi]\in\WFF).
\]
\end{remark*}

\begin{remark*}[Definition predicate reading]
\[
\operatorname{FormulaContext}(C[-])
\]
means that \(C[-]\) is a one-hole propositional formula context.
\end{remark*}

\begin{remark*}[Interpretation]
A formula context is a syntactic environment into which a formula can be inserted. Contexts make precise the idea of replacing a subformula inside a larger formula.
\end{remark*}

\begin{remark*}[Examples]
If
\[
C[-] := ([-]\land R),
\]
then
\[
C[P]=(P\land R)
\]
and
\[
C[Q\lor S]=((Q\lor S)\land R).
\]
\end{remark*}

\begin{dependencies}
\begin{itemize}
  \item \hyperref[def:well-formed-formula]{Well-Formed Formula}
  \item \hyperref[def:subformula-propositional-logic]{Subformula}
\end{itemize}
\end{dependencies}
% END GENERATED ARTIFACT: def:formula-context-propositional-logic
% BEGIN GENERATED ARTIFACT: thm:replacement-of-logical-equivalents-propositional-logic
\begin{theorembox}{Theorem (Replacement of Logical Equivalents)}
\begin{theorem}[Replacement of Logical Equivalents]
\label{thm:replacement-of-logical-equivalents-propositional-logic}
Replacing logically equivalent formulas inside a formula context preserves logical equivalence:
\[
\varphi\equiv\psi
\Longrightarrow
C[\varphi]\equiv C[\psi].
\]

\noindent\hyperref[prf:replacement-of-logical-equivalents-propositional-logic]{Go to proof.}
\end{theorem}
\end{theorembox}

\begin{remark*}[Standard quantified statement]
\[
(\forall\varphi,\psi\in\WFF)
(\forall C[-])
\left[
\operatorname{FormulaContext}(C[-])
\land
\varphi\equiv\psi
\to
C[\varphi]\equiv C[\psi]
\right].
\]
\end{remark*}

\begin{remark*}[Predicate reading]
\[
\varphi\equiv\psi
\quad\Longrightarrow\quad
C[\varphi]\equiv C[\psi].
\]
A logical equivalence may be used inside a larger formula context without changing
the truth behavior of the whole formula.
\end{remark*}

\begin{remark*}[Interpretation]
Logical equivalence is substitutive: equivalent formulas may be replaced inside larger formulas without changing truth behavior.
\end{remark*}

\begin{remark*}[Exposition]
The equivalence laws in this section may be used as replacement rules inside
larger formulas. If a formula contains a subformula matching one side of a listed
equivalence law, that subformula may be replaced by the other side.

Common replacement rules include:
\[
\neg\neg\varphi \equiv \varphi,
\]
\[
\neg(\varphi\land\psi)\equiv \neg\varphi\lor\neg\psi,
\qquad
\neg(\varphi\lor\psi)\equiv \neg\varphi\land\neg\psi,
\]
\[
\varphi\to\psi\equiv \neg\varphi\lor\psi,
\]
\[
\varphi\leftrightarrow\psi
\equiv
(\varphi\to\psi)\land(\psi\to\varphi),
\]
and
\[
\varphi\leftrightarrow\psi
\equiv
(\varphi\land\psi)\lor(\neg\varphi\land\neg\psi).
\]
\end{remark*}

\begin{remark*}[Exposition]
The laws in this section allow formulas to be manipulated algebraically. They are not new formation rules; they are semantic equivalences. Later sections use these laws to transform formulas into normal forms and to reduce the set of primitive connectives.
\end{remark*}

\begin{remark*}[Exposition]
The algebra of propositions is the bridge from semantics to computation: once formulas are equivalent, they can be rewritten without changing truth behavior.
\end{remark*}

\begin{dependencies}
\begin{itemize}
  \item \hyperref[def:formula-context-propositional-logic]{Formula Context}
  \item \hyperref[def:logical-equivalence-propositional-logic]{Logical Equivalence}
  \item \hyperref[lem:coincidence-lemma-propositional-logic]{Coincidence Lemma for Propositional Logic}
  \item \hyperref[thm:unique-extension-truth-assignment-propositional-logic]{Unique Extension of a Truth Assignment}
\end{itemize}
\end{dependencies}
% END GENERATED ARTIFACT: thm:replacement-of-logical-equivalents-propositional-logic
